\chapter[Dictionary and misfit]{The universal dictionary and the misfit certificate}\label{ch:cIII5}
\chapterbrief
{$\Omega\subset\R^N$ open, $N=1+d$, $R=\C[\tau,\xi]$; a \emph{filtered dictionary} $\mathcal D_0\subset\mathcal D_1\subset\cdots$ of finite sets of operators (atoms), $\mathcal D_\infty=\bigcup\mathcal D_n$; a \emph{level} $n$ is $\mathrm{span}\,\mathcal D_n=\mathrm{span}\{D_0,\dots,D_{N_n}\}$; $u$ \emph{fits} at level $n$ if some non-zero $c$ gives $\sum c_iD_iu=0$, otherwise it has a \emph{misfit}.}
{to prove what ``a dictionary that takes everything into account'' can and cannot be: it exists in the sense of approximation (Peetre's theorem), it cannot give exact identification of every operator (Baire), a misfit can be certified by finite data but a fit never can, and a stopping rule exists with a guarantee under a threshold that the data cannot supply.}
{Dim.~A48--A51. They prove Chapter~\ref{PAP-ch:dictionary} of the paper.}

\section{Universality}

\begin{dimcard}{A48}{DICT-U1}{Every local linear operator is a limit of dictionary elements}{\proved}{thm:dictu1}
\textbf{(U1-a)} Let $T:C^\infty(\Omega)\to C^\infty(\Omega)$ be linear with $\mathrm{supp}\,Tu\subset\mathrm{supp}\,u$. Then $Tu=\sum_\alpha a_\alpha\partial^\alpha u$ with $a_\alpha\in C^\infty(\Omega)$ locally finite (on each compact only $|\alpha|\le r(K)$ occur).

\textbf{(U1)(a)} Let $(\chi_j)$ be an explicit locally finite smooth partition of unity subordinate to an exhaustion of $\Omega$, and $\mathcal D^{\rm loc}_n=\{\chi_j(y)y^\beta\partial^\alpha:\ j,|\alpha|,|\beta|\le n\}$. Every linear local $T$ is a limit, in the topology $\tau_1$ (locally bounded order, coefficients converging in $C^\infty(K)$), of elements of $\mathrm{span}_{\mathbb Q(i)}\mathcal D^{\rm loc}_\infty$.

\textbf{(N0)} With $\mathcal D^{\rm pol}_n=\{y^\beta\partial^\alpha:|\alpha|,|\beta|\le n\}$, $T$ is a pointwise limit of operators with polynomial coefficients iff $T$ has globally finite order on $\Omega$.
\end{dimcard}

\begin{proofidea}
Peetre's theorem says a linear local operator has locally finite order; then the coefficients are read off by applying $T$ to monomials, and they are approximated by polynomials. N0 uses Banach--Steinhaus: equicontinuity forces a uniform order.
\end{proofidea}

\begin{proof}
\step{1}{(U1-a) Peetre's theorem (cited)}
If $D:C^\infty\to C^\infty$ is linear with $\mathrm{supp}\,Ds\subset\mathrm{supp}\,s$, then for each compact $K$ there is $r$ with $j^r_ys=0\Rightarrow Ds(y)=0$ for $y\in K$ (Peetre 1959--60; Kol\'a\v r--Michor--Slov\'ak, Thm 19.1; \cf). So $j^r_yu\mapsto Tu(y)$ is well defined and linear on jets, hence $Tu(y)=\sum_{|\alpha|\le r}a_\alpha(y)\partial^\alpha u(y)$. For smoothness of $a_\alpha$, apply $T$ to $m_\beta=\chi z^\beta/\beta!$ near $y$: $T(m_\beta)(y)=\sum_{\alpha\le\beta}a_\alpha(y)y^{\beta-\alpha}/(\beta-\alpha)!$ is a triangular system for the $a_\alpha$ with smooth data.

\step{2}{(U1)(a) Approximation}
Polynomials are dense in $C^\ell(K)$: for $g$ smooth near $K$, mollify $\chi g$ with a Gaussian to get an entire function, whose Taylor series converges with all derivatives uniformly on $K$, then round the coefficients to $\mathbb Q(i)$. Let $r_j$ be the order of $T$ on $\mathrm{supp}\,\chi_j$; choose polynomials $p_{\alpha,j,n}$ with $\|a_\alpha-p_{\alpha,j,n}\|_{C^n(\mathrm{supp}\chi_j)}\le2^{-n}/(1+\|\chi_j\|_{C^n})$ and put $T_n=\sum_{j\le n}\chi_j\sum_{|\alpha|\le r_j}p_{\alpha,j,n}\partial^\alpha$. On each compact only finitely many $\chi_j$ matter, and $T_n\to T$ in $\tau_1$.

\step{3}{(N0) Pointwise limits}
($\Leftarrow$) If the order is $\le k$ globally, choose polynomials $p_{\alpha,n}$ with $\|a_\alpha-p_{\alpha,n}\|_{C^n(K_n)}\le1/n$. ($\Rightarrow$) $C^\infty(\Omega)$ is Fr\'echet, so a pointwise convergent sequence $\{T_n\}$ is equicontinuous (Banach--Steinhaus): on a compact $K$ with interior there are $k,K',C$ with $\sup_K|T_nu|\le C\|u\|_{C^k(K')}$ for all $n,u$. A $T_n$ with a non-zero coefficient of order $>k$ would contradict this bound; the contradiction is obtained from the principal symbol of the highest-order coefficient at a point of $\mathrm{int}\,K$ (\emph{condensed}: the explicit test functions are in the package \S1.1).

\whyworks{a local operator sees a function only through finitely many derivatives at each point (Peetre), so it is a differential operator with smooth coefficients, and smooth coefficients are polynomial limits.}
\end{proof}
\uses{Peetre (\cf), Banach--Steinhaus.}

\begin{example}{An operator of infinite global order}
On $\Omega=\R$ take $T=\sum_k\varphi_k(x)\partial_x^k$ with $\varphi_k\in C_c^\infty((k,k+1))$ non-zero. It is local and linear, of globally infinite order: it is not a pointwise limit of polynomial-coefficient operators, but it is a limit in $\mathcal D^{\rm loc}$ by (a).
\end{example}

\section{The negative theorem}

\begin{dimcard}{A49}{DICT-N1}{No countable dictionary identifies everything}{\proved}{thm:dictn1}
\begin{enumerate}[label=(\roman*)]
\item (Baire) Let $X$ be an infinite-dimensional $F$-space and $\{a_k:P_k\to X\}$ countably many continuous maps from $\sigma$-compact $P_k\subset\R^{p_k}$. Then $\bigcup_ka_k(P_k)$ is meagre. In particular for any countable dictionary the exactly representable operators form a meagre set in each class of coefficients.
\item No finite set of data certifies the fit or the identification (A50, A7, A9).
\item (iii) (Non-monotonicity) If $N\ge2$, $u\in\EP(\Omega)$, $u\not\equiv0$, $D=\dim R\cdot u$, then for every $n$ with $\binom{n+N-1}{N-1}>D$ and every $L\in\Ann(u)$ with $\deg L\ge1$: $\Ann(u)\cap R_{\le n}\supsetneq(L)\cap R_{\le n}$.
\item (iv) A procedure with finite outputs over a finite alphabet (also as a stabilising limit) emits only countably many operators, so it never emits $\partial_t-c\partial_x$ for all but countably many $c$.
\item (v) Let $V_0\subset V_1\subset\cdots$ be finite-dimensional spaces of differential polynomials in $u$ and $B_n=\{u:\exists P\in V_n\setminus0,\ P[u]=0\}$. Then $B_n$ is closed and nowhere dense (so $\bigcup B_n$ is meagre) and is \emph{shy} (Christensen): the generic $u$ satisfies no equation of the dictionary, in the topological sense and in the sense of measure.
\end{enumerate}
\end{dimcard}

\begin{proofidea}
(i) A compact set in an infinite-dimensional space has empty interior (Riesz), so countably many images of compacts are meagre. (v) Move $u$ in a ``direction of escape'' $w$ along which the determinant of the dictionary applied to $u+sw$ is a non-zero polynomial in $s$, so only finitely many $s$ lie in $B_n$.
\end{proofidea}

\begin{proof}
\step{1}{(i) Baire}
Each $a_k(P_k)$ is a countable union of images of compacts, i.e.\ of compacts. A compact $C$ in a Hausdorff topological vector space of infinite dimension has empty interior: if $C\supset x+V$ with $V$ a neighbourhood of $0$, $V$ would lie in a compact neighbourhood of $0$, and a locally compact Hausdorff TVS has finite dimension (Riesz). $C$ is closed, hence nowhere dense; apply Baire. (With rational parameters each $a_k$ has countable image, and the conclusion holds without any continuity.)

\step{2}{(iii) A dimension count}
$Q\mapsto Qu$ on $R_{\le n}$ has kernel $\Ann(u)\cap R_{\le n}$ and image of dimension $\le D$, so $\dim(\Ann(u)\cap R_{\le n})\ge\binom{n+N}N-D$. Also $\dim((L)\cap R_{\le n})=\dim R_{\le n-\deg L}\le\binom{n-1+N}N$. Since $\binom{n+N}N-\binom{n+N-1}N=\binom{n+N-1}{N-1}>D$, the first dimension is larger, and $(L)\cap R_{\le n}\subset\Ann(u)\cap R_{\le n}$ because $L\in\Ann(u)$. (For $N=1$ the statement is false: $\Ann(u)$ is principal.)

\step{3}{(iv) Countability}
Finite strings over a finite alphabet are countable; a stabilised sequence has a limit string; the image of a countable set is countable. The family $\{\partial_t-c\partial_x\}_{c\in\R}$ is uncountable and each member is identified by suitable solutions (C3-T4).

\step{4}{(v) Escape directions}
Take a degree-adapted basis $P_1,\dots,P_M$ of $V_n$ (top-degree parts $P_i^{\rm top}$ linearly independent; possible by choosing complements in the filtration by degree), $e_i=\deg P_i$. \emph{Closed:} if $u_k\to u$ with $\sum c^{(k)}_iP_i[u_k]=0$, $|c^{(k)}|=1$, pass to a convergent subsequence $c^{(k)}\to c$ and use continuity of $P_i[\cdot]$. \emph{Escape:} by the non-degeneracy lemma (ND-NL: if $Q_1,\dots,Q_M$ are linearly independent as functions on $U\times J^r$, there is $w\in C_c^\infty(U)$ with $Q_i[w]$ independent) take $w$ with $P_1^{\rm top}[w],\dots,P_M^{\rm top}[w]$ independent and test functionals $\ell_j$ with $\det(\ell_j(P_i^{\rm top}[w]))\ne0$. For any $u$, $P_i[u+sw]=s^{e_i}P_i^{\rm top}[w]+(\text{lower in }s)$, so $p_u(s)=\det(\ell_j(P_i[u+sw]))$ is a polynomial in $s$ with non-zero coefficient of $s^{\sum e_i}$. If $u+sw\in B_n$ the $P_i[u+sw]$ are dependent and $p_u(s)=0$: at most finitely many $s$. So the complement of $B_n$ is dense (meagre), and the probability measure of $sw$, $s\in[0,1]$ uniform, gives $\mu(B_n-u)=0$ for every $u$: $B_n$ is shy; a countable union of shy sets is shy.

\whyworks{a countable list of finite-dimensional families can only cover a ``thin'' set, and the thinness is of both kinds, topological and measure-theoretic, because every direction of escape is a polynomial identity in one real parameter.}
\end{proof}
\uses{A2, A7, A9, A50.}

\begin{caution}[title=What the theorem does not say]
``Everything'' holds in the sense of approximation (A48) plus a misfit certificate (A50), not as exact identification. For a differential operator with $C^\infty$ coefficients and one data function the statement ``no operator is identified'' (the first form of N1(vi)) is false; it was withdrawn and replaced by the rank theorem (Chapter~\ref{ch:cIV3}).
\end{caution}

\section{Fit and misfit}

\begin{dimcard}{A50}{DICT-A}{Misfit has a finite certificate; fit never does (Theorem A, MF1--MF3)}{\proved}{thm:dicta}
Let $\mathrm{span}\{D_0,\dots,D_N\}$ be a level of linear continuous operators $C^\infty(\Omega)\to\mathcal D'(\Omega)$, independent over $\C$; $U_n=\{u:D_0u,\dots,D_Nu\text{ independent}\}$ (misfit), $F_n=C^\infty\setminus U_n$ (fit).
\begin{enumerate}[label=(\roman*)]
\item $U_n$ is open, and every $u\in U_n$ has a finite certificate ($N+1$ test functions) valid in a neighbourhood.
\item $F_n$ is closed with empty interior: no finite set of \emph{approximate} data certifies a fit.
\item For differential atoms and any continuous functionals $\ell_1,\dots,\ell_M$ on $C^\infty(\Omega)$, every $u\in F_n$ has a $v\in U_n$, arbitrarily close, with $\ell_j(v)=\ell_j(u)$: a fit is not certifiable even from \emph{exact} finite data.
\end{enumerate}
Certificates: (MF1) $\mathrm{dist}(v,W)=\max\{|\varphi(v)|:\|\varphi\|\le1,\varphi|_W=0\}$; (MF2) in a Hilbert space $r_n^2=\det G(D_0u,\dots,D_Nu)/\det G(D_1u,\dots,D_Nu)$; (MF3) weak form: independence is equivalent to $\det\langle f_i,\psi_j\rangle\ne0$ for some test functions from a fixed countable dense set, with robustness $|\sigma_k(A)-\sigma_k(B)|\le\|A-B\|$.
\end{dimcard}

\begin{proofidea}
Linear independence is an open condition witnessed by a non-vanishing determinant, so misfit is open and certified by finitely many test functions. A fit is a closed condition with empty interior because locality lets one perturb $u$ on disjoint small balls to restore independence without changing finitely many functionals.
\end{proofidea}

\begin{proof}
\step{1}{(MF1) The distance formula}
For $w\in W$, $|\varphi(v)|=|\varphi(v-w)|\le\|v-w\|$. Conversely with $\delta=\mathrm{dist}(v,W)>0$ define $\varphi_0(w+\lambda v)=\lambda\delta$ on $\mathrm{span}(W,v)$; since $\|\lambda v+w\|\ge|\lambda|\delta$, $\|\varphi_0\|\le1$, and Hahn--Banach extends it with $\varphi(v)=\delta$.

\step{2}{(MF2) The Gram quotient}
Replacing the first row and column of $G(D_0u,\dots)$ by those of $\mathrm{res}=D_0u-P_WD_0u$ (elementary operations) gives $\mathrm{diag}(\|\mathrm{res}\|^2,G(D_1u,\dots))$, since $\mathrm{res}\perp W$; and $\varphi=\langle\cdot,\mathrm{res}\rangle/\|\mathrm{res}\|$ has $\varphi(D_iu)=0$, $\varphi(D_0u)=\|\mathrm{res}\|$.

\step{3}{(MF3) The weak form}
$f_0,\dots,f_N\in\mathcal D'$ are independent iff some $\det\langle f_i,\psi_j\rangle\ne0$: if $\sum c_if_i=0$ the columns are dependent; otherwise the image of $\psi\mapsto(\langle f_i,\psi\rangle)_i$ is a proper subspace orthogonal to some $c\ne0$, giving $\sum\bar c_if_i=0$. For differential atoms $\langle D_iu,\psi\rangle=\langle u,D_i^t\psi\rangle$: no derivative of $u$ is needed.

\step{4}{(i)--(ii)}
(i) By MF3, $\det A(u)\ne0$ for some $\psi_j$; $u\mapsto A(u)$ is continuous on $C^\infty$ (and in the $\mathcal D'$ topology when $D_i^t\psi_j\in C_c^\infty$), so $\det A(v)\ne0$ nearby. (ii) $F_n$ is the complement of (i); it has empty interior by A49(v) with the non-degeneracy lemma, and a set of approximate data contains a neighbourhood of $u$.

\step{5}{(iii) Exact finite data}
Let $Z=\bigcap_j\ker\ell_j$, closed of codimension $\le M$. The idea is to use locality: perturb $u$ by a function $w\in Z$ built from pieces supported in small disjoint balls, so that the perturbation does not change the finitely many data but restores independence of the $D_iv$, and $v=u+\varepsilon w$ lies in $U_n$ for small $\varepsilon\ne0$. (\emph{Condensed}: the choice of the balls is in the package \S3.3, where a flaw of the first version, ``independence on an open set implies independence on each ball'', was repaired; numerical check: $u=e^{-t}\sin x$ with fit in $\{\partial_t,\partial_x^2\}$, three functionals, relative $\lambda_{\min}=4.5\cdot10^{-2}$.) The statement is open for non-local atoms.

\whyworks{the asymmetry is topological: ``$\ne$'' is open, ``$=$'' is closed and thin; and locality lets one repair equality by a perturbation invisible to finitely many measurements.}
\end{proof}
\uses{A49.}

\begin{example}{Certifying a misfit}
Test whether $u=e^{-t}\sin x$ satisfies a law in the level $\{\partial_t,\partial_x\}$: the Gram matrix of $(u_t,u_x)$ in $L^2$ is non-singular with $\lambda_{\min}>0$, so $\varphi=\langle\cdot,\mathrm{res}\rangle/\|\mathrm{res}\|$ is a finite certificate that no first-order law holds. The same $u$ \emph{does} fit the level $\{\partial_t,\partial_x^2\}$ ($u_t=-u$, $u_{xx}=-u$), and no finite data can certify that.
\end{example}

\begin{dimcard}{A51}{DICT-GR}{Stopping rule (G-R) and identification in the limit with a dictionary (I-LIM)}{\proved}{thm:dictgr}
\textbf{(G-R)} Let $n_0$ be the least level with an affine fit, $L\in\mathcal A_{n_0}$ with $Lu=0$, and assume (H-ID): $\Ann(u)\cap\mathrm{span}\,\mathcal D_n=\langle L\rangle_n$ for all $n\ge n_0$. Then the test R1 ($r_n=\mathrm{dist}(D_0u,W_n)$) stops exactly at $n_0$, and the rule R2 (coherence) signals failure when the kernel at a higher level is strictly larger than $\langle L\rangle_n$.

\textbf{(I-LIM, approximate branch)} Let $\Omega$ be bounded Lipschitz, $u\in H^m$, $\mathcal K_R=\{L=D_0+\sum_{\alpha\in A}c_\alpha\partial^\alpha:\|c_\alpha\|_{C^{s,1}}\le R\}$ (compact in $C^s$), (H-ID) uniqueness in $\mathcal K_{R'}$ and (H2) bounded approximability by $\tilde L_n\in\mathrm{span}\,\mathcal D_n\cap\mathcal K_{R'}$ with $\|\tilde L_n-L\|_{C^s}\le\varepsilon_n\to0$. Then the modulus $\omega(\eta)=\sup\{\|L'-L\|_{C^s}:L'\in\mathcal K_{R'},\ \|L'u\|_{L^2}\le\eta\}\to0$, every minimiser $L_n$ of $\|L'u\|$ on $\mathrm{span}\,\mathcal D_n\cap\mathcal K_{R'}$ satisfies $\|L_n-L\|\le\omega(C_u\varepsilon_n)$, $C_u=\sum_\alpha\|\partial^\alpha u\|_{L^2}$, and with data $\tilde u$, $\|\tilde u-u\|_{H^m}\le\delta$: $\|L_n^\delta-L\|\le\omega(C_u\varepsilon_n+2C_{R'}\delta)$.
\end{dimcard}

\begin{proofidea}
Compactness of $\mathcal K_{R'}$ plus uniqueness (H-ID) gives that small residual means close coefficients; the rest is the triangle inequality.
\end{proofidea}

\begin{proof}
\step{1}{(a) The modulus tends to $0$}
Otherwise there are $\varepsilon_0>0$ and $L'_j\in\mathcal K_{R'}$ with $\|L'_ju\|\to0$, $\|L'_j-L\|\ge\varepsilon_0$. By compactness $L'_j\to L^*$ in $C^s$ with $L^*\in\mathcal K_{R'}$; the map (coefficients)$\mapsto L'u$ is continuous $C^0\to L^2$ ($\|(L'-L'')u\|\le\sum\|c'_\alpha-c''_\alpha\|_\infty\|\partial^\alpha u\|_{L^2}$), so $L^*u=0$ and $\|L^*-L\|\ge\varepsilon_0$, against (H-ID).

\step{2}{(b) and (c)}
$\|L_nu\|\le\|\tilde L_nu\|=\|(\tilde L_n-L)u\|\le C_u\varepsilon_n$, then apply (a). For data $\tilde u$ and $L'\in\mathcal K_{R'}$, $\|L'\tilde u-L'u\|\le C_{R'}\delta$ with $C_{R'}=\|D_0\|_{H^m\to L^2}+|A|R'$; hence $\|L_n^\delta u\|\le\|L_n^\delta\tilde u\|+C_{R'}\delta\le\|\tilde L_n\tilde u\|+C_{R'}\delta\le C_u\varepsilon_n+2C_{R'}\delta$.

\step{3}{(G-R) The stopping test}
By (H-ID), for $n<n_0$ there is no affine fit, so $r_n>0$, and at $n_0$ the fit exists, so $r_{n_0}=0$ (computed by MF2 on a maximal independent subfamily of $W_n$). For $n>n_0$ the coherence rule R2 compares $\Ann(u)\cap\mathrm{span}\,\mathcal D_n$ with $\langle L\rangle_n$; a strictly larger kernel means $u$ is a closed-form datum or does not identify (A49(iii), A5).

\step{4}{When (H-ID) holds}
With a single coefficient and $\{\partial^{\alpha_1}u\ne0\}$ dense, (H-ID) holds for every $R'$ (rank theorem, Chapter~\ref{ch:cIV3}). For $|A|=1$, $D_0=\partial_t$, $|\partial^{\alpha_1}u|\ge\gamma>0$ and a cone condition on $\Omega$, the modulus is H\"older of exponent $2/(N+2)$: $\|L^\delta_{n(\delta)}-L\|_{C^0}=O((\varepsilon_n+\delta)^{2/(N+2)})$.

\whyworks{identification in the limit with a dictionary is a compactness argument; what the data can never supply is the margin: the rate depends on $\varepsilon_n$, the approximability of the true law.}
\end{proof}
\uses{A41, A43, A49, A50.}
